% !TeX program = pdflatex
% !TeX encoding = UTF-8
%
% Унифицированный отказоустойчивый адаптивный контроллер UFTC:
% открытомировое обнаружение отказов и иерархическая инкрементальная адаптация
% летательного аппарата.
%
% Подготовлено для подачи в журнал «Труды МАИ» (ВАК).
%
\documentclass[a4paper,14pt]{extarticle}

\usepackage[T2A]{fontenc}
\usepackage[utf8]{inputenc}
\usepackage[english,russian]{babel}

\usepackage[a4paper,top=20mm,bottom=20mm,left=25mm,right=15mm]{geometry}
\usepackage{setspace}
\onehalfspacing

\usepackage{amsmath,amssymb,amsthm}
\newtheorem{lemma}{Лемма}
\usepackage{mathtools}
\usepackage{bm}
\usepackage{graphicx}
\usepackage{caption}
\captionsetup{font=small,labelfont=bf,labelsep=endash}
\usepackage{booktabs}
\usepackage{array}
\usepackage{tabularx}
\usepackage{indentfirst}
\usepackage{microtype}
\usepackage{enumitem}
\setlist{leftmargin=*,nosep}

\usepackage[unicode,colorlinks=true,linkcolor=black,citecolor=black,urlcolor=blue]{hyperref}

\usepackage{tikz}
\usetikzlibrary{positioning,arrows.meta,shapes.geometric,calc,fit,backgrounds}
\usepackage{float}
\newcounter{algorithm}

% Сокращения
\newcommand{\UFTC}{\textsc{UFTC}}
\newcommand{\ET}{\textsc{ET-DHP}}
\newcommand{\IADP}{\textsc{iADP}}
\newcommand{\AINDI}{\textsc{A-INDI}}
\newcommand{\INDI}{\textsc{INDI}}
\newcommand{\FDD}{\textsc{FDD}}
\newcommand{\CUSUM}{\textsc{CUSUM}}
\newcommand{\RLS}{\textsc{RLS}}
\newcommand{\degr}{^{\circ}}

% Русские обозначения рулей
\newcommand{\dv}{\delta_{\text{в}}}
\newcommand{\dn}{\delta_{\text{н}}}
\newcommand{\de}{\delta_{\text{э}}}
\newcommand{\dt}{\delta_{\text{т}}}

\title{}
\date{}

\begin{document}

\thispagestyle{empty}

\begin{flushleft}
\small УДК~629.7.072.1:629.7.05:519.876.5
\end{flushleft}

\begin{center}
\large\bfseries
Унифицированный отказоустойчивый адаптивный контроллер UFTC с обнаружителем отказов открытого мира и иерархической инкрементальной адаптацией
\end{center}

\medskip

\begin{center}
\textbf{Мазаев~А.\,С.\textsuperscript{1,*}}
\end{center}

\begin{center}
\small
\textsuperscript{1}\,Московский авиационный институт (национальный исследовательский университет), \\
Волоколамское шоссе, 4, Москва, 125993, Россия \\[2pt]
\textsuperscript{*}\,e-mail: simplex.bro@gmail.com
\end{center}

\bigskip

% --------------------------------------------------------------------
% Аннотация
% --------------------------------------------------------------------
\noindent\textbf{Аннотация.}
\emph{Цель работы}~--- разработка отказоустойчивого адаптивного контроллера летательного аппарата, ориентированного на парирование заранее не каталогизированных аномалий динамики без предобученных наблюдателей под каждый режим отказа. \emph{Подход:} предложена архитектура унифицированного отказоустойчивого адаптивного контроллера \UFTC{} (Unified Fault-Tolerant Control)~--- иерархическая композиция инкрементального приближённого динамического программирования (\IADP{}) на среднем уровне, активно-адаптивного инкрементального обращения нелинейной динамики (\AINDI{}) на внутреннем уровне и параллельного обнаружителя аномалий, использующего единственный номинальный фильтр Калмана с адаптацией ковариаций по схеме Сейджа--Хузы и тест разладки CUSUM по нормированной махаланобисовой дистанции инноваций. Восстановление контроллера после обнаружения реализовано через инфляцию ковариационной матрицы рекурсивных наименьших квадратов в \IADP{} и временное снижение коэффициента забывания с последующим линейным возвратом к номинальному значению; связь между внутренним и средним уровнями обеспечивается непрерывной модуляцией радиуса доверительной области по сигналу тяжести отказа. \emph{Теоретический результат:} получена корректировка формул адаптации Сейджа--Хузы (Лу, 2015) с проекцией собственных значений на относительный нижний предел, гарантирующая положительную определённость матриц $Q$ и $R$ при произвольных начальных условиях. \emph{Численный эксперимент:} реализация \UFTC{} в составе библиотеки \texttt{TensorAeroSpace} верифицирована 56 модульными тестами (покрытие 93\,\%) и интеграционными прогонами на нелинейной угловой модели F-16 (шесть пресетов повреждений) и на нелинейной 6-DoF модели Boeing~747. В 300-секундном сценарии отказа крайнего левого двигателя B-747 полносостояниная версия \UFTC{} использует вектор $[e_V,e_h,e_{\vartheta},e_{\psi},e_{\varphi},e_{\beta},p,q,r]$ без бинарного признака отказа; \FDD{} срабатывает на $t=10{,}05$~с, а RMSE в последних 20~с составляет $0{,}0010$~ft/s по скорости, $0{,}0012$~ft по высоте и $4\cdot10^{-6}\,\degr/0{,}000620\degr/0{,}000092\degr$ по тангажу, курсу и крену. Разомкнутый контур к $t=300$~с уходит на $\Delta h=-61\,392$~ft и $\Delta\psi=-43{,}04\degr$.

\medskip
\noindent\textbf{Ключевые слова:}
отказоустойчивое управление; адаптивное управление; инкрементальное динамическое программирование; INDI; обнаружение отказов; CUSUM; нелинейная динамика самолёта; Boeing~747; F-16; открытый мир.

\bigskip

% --------------------------------------------------------------------
% English abstract
% --------------------------------------------------------------------
\selectlanguage{english}
\noindent\textbf{Unified fault-tolerant adaptive controller UFTC with an open-world fault detector and hierarchical incremental adaptation}

\smallskip
\noindent\textbf{A.\,S.~Mazaev\textsuperscript{1,*}}

\smallskip
\noindent
\small
\textsuperscript{1}\,Moscow Aviation Institute (National Research University), Volokolamskoe sh.\ 4, Moscow, 125993, Russia \\
\textsuperscript{*}\,e-mail: simplex.bro@gmail.com

\medskip
\noindent\textbf{Abstract.}
The paper introduces \UFTC{} (Unified Fault-Tolerant Control), a hierarchical adaptive flight-control architecture composing incremental Approximate Dynamic Programming (\IADP{}) on the middle layer, Active-Adaptive Incremental Nonlinear Dynamic Inversion (\AINDI{}) on the inner layer, and a parallel anomaly detector that requires no a-priori fault catalog. Unlike multimodel adaptive estimation (MMAE) banks, \UFTC{} uses a single nominal Sage--Husa Kalman filter with a CUSUM change-point test on the normalized Mahalanobis distance of innovations, which is aimed at detecting deviations from nominal dynamics without pre-trained mode-specific observers. Post-detection recovery is unified through covariance inflation of the recursive least-squares inside \IADP{} and a temporary forgetting-factor drop with linear ramp back to nominal. The coupling between the inner and the middle layers is realized through a continuously modulated trust-region radius driven by the detector severity signal. We derive a correction to Lu 2015 Sage--Husa adaptation formulas using eigenvalue projection onto a relative floor, which guarantees positive-definiteness of $Q$ and $R$ under arbitrary initial conditions. The \UFTC{} implementation in the \texttt{TensorAeroSpace} library is verified by 56 unit tests with 93\,\% coverage and integration runs on the nonlinear angular F-16 model with six damage presets and on the nonlinear 6-DoF Boeing~747 model. In the 300-s B-747 left-outer-engine failure scenario, the full-state \UFTC{} uses $[e_V,e_h,e_{\theta},e_{\psi},e_{\phi},e_{\beta},p,q,r]$ without a binary fault indicator; the FDD alarm fires at $t=10.05$~s, and the last-20-s RMSE is $0.0010$~ft/s in speed, $0.0012$~ft in altitude, and $4\cdot10^{-6}$/$0.000620$/$0.000092$~deg in pitch/yaw/roll. The open loop reaches $\Delta h=-61\,392$~ft and $\Delta\psi=-43.04$~deg by $t=300$~s.

\medskip
\noindent\textbf{Keywords:} fault-tolerant control; adaptive control; incremental dynamic programming; INDI; fault detection; CUSUM; nonlinear aircraft dynamics; Boeing~747; F-16; open world.
\selectlanguage{russian}

\bigskip

% ====================================================================
\section{Введение}
% ====================================================================

\paragraph{Терминологические соглашения.} В тексте используются следующие сокращения и обозначения, согласованные с традицией российской школы автоматического управления летательными аппаратами~\cite{borodin1978}: ОУ~--- объект управления (нелинейная 6-DoF модель самолёта); САУ~--- система автоматического управления; УУ~--- управляющее устройство (регулятор); ОС~--- обратная связь; FTC (Fault-Tolerant Control)~--- отказоустойчивое управление; FDD (Fault Detection and Diagnosis)~--- модуль обнаружения и диагностики отказов; $u\in\mathbb{R}^{n_u}$~--- вектор управляющих воздействий, $x\in\mathbb{R}^{n_x}$~--- вектор состояния ОУ; $\Delta(\cdot)$~--- вариация (отклонение) от невозмущённого режима по терминологии метода малых возмущений. Обозначения рулевых поверхностей следуют~\cite{borodin1978}: $\dv$~--- руль высоты, $\de$~--- элероны, $\dn$~--- руль направления, $\dt$~--- орган управления тягой (РУД).

С точки зрения классической теории автоматического управления рассматриваемая в работе задача формулируется как стабилизация нелинейного многосвязного ОУ при структурных изменениях оператора объекта. Изменения возникают вследствие отказов исполнительных механизмов (двигателей, рулевых поверхностей) или повреждений планера. Качество замкнутой системы оценивается стандартным набором показателей: ошибка стабилизации заданного режима полёта, время восстановления управляемости, апериодичность переходного процесса по угловым скоростям, отсутствие установившейся ошибки. Принципиальное отличие настоящей задачи от классической задачи стабилизации~--- в \emph{неполноте априорной информации об операторе ОУ}: оператор изменяется в произвольный момент времени по неизвестному заранее сценарию, что исключает применение фиксированной структуры регулятора. Эту особенность снимают адаптивные методы синтеза с онлайн-идентификацией параметров.

Современные требования к безопасности коммерческих и беспилотных летательных аппаратов делают отказоустойчивое управление (Fault-Tolerant Control, FTC) одной из центральных задач теории и практики авиационного управления. Классическая постановка опирается на схему \emph{обнаружение → классификация → переключение}: модуль обнаружения и идентификации отказа (Fault Detection and Identification, FDI) определяет текущий режим отказа из заранее заданного каталога, после чего управляющий контур переключается на предобученный закон, специфический для этого режима. Многомодельное адаптивное оценивание (MMAE)~\cite{lu2015,lu2017} даёт каноническую реализацию такой схемы: $K+1$ параллельных фильтров Калмана, каждый настроенный под конкретный режим отказа, формируют апостериорное распределение по режимам, и контроллер выбирает соответствующий закон. Распространение событийно-триггерных подходов адаптивной критики~\cite{sun2022stab,konatala2024} позволило вынести часть онлайн-обучения в полётный цикл, однако сохранилось ключевое ограничение: вектор состояния агента расширяется явным признаком отказа $\ell$, а закон управления обучается \emph{закрытомиром} на конкретный сценарий повреждения, заданный заранее.

Допущение о закрытом каталоге отказов ограничивает применимость FTC к реальным эксплуатационным условиям, в которых полный набор сценариев априорно неизвестен. Боевые повреждения, попадание птиц, частичные потери эффективности рулевых поверхностей и комбинированные сенсорно-актуаторные отказы могут лежать вне обучающего распределения. Альтернативная парадигма~--- управление открытого мира (open-world FTC) с непрерывной онлайн-адаптацией к произвольной аномалии динамики~--- развивается в направлении инкрементальных нелинейных подходов на основе обращения динамики~\cite{sieberling2010,wang2019}, активно-адаптивного INDI~\cite{aaindi2024} и инкрементального ADP~\cite{zhou2018,konatala2024}. Каждый из этих подходов решает часть проблемы: \AINDI{} обеспечивает робастность внутреннего контура к изменениям эффективности управляющих поверхностей, \IADP{} компенсирует постепенный дрейф динамики через онлайн-идентификацию локальной модели, однако единая архитектура, объединяющая \emph{обнаружение} и \emph{адаптацию} без априорного каталога, в литературе отсутствует.

Настоящая работа предлагает такую архитектуру~--- унифицированный отказоустойчивый адаптивный контроллер \UFTC{} (Unified Fault-Tolerant Control), в основе которого лежат следующие проектные решения:

\begin{enumerate}
\item Иерархическая композиция трёх слоёв (внутренний \AINDI{}, средний \IADP{}, параллельный обнаружитель аномалий) с явно определёнными интерфейсами и согласованным масштабом адаптации.
\item Обнаружение в режиме открытого мира через тест разладки \CUSUM{} по нормированной махаланобисовой дистанции инноваций единственного номинального фильтра Калмана~--- без предобучения банка наблюдателей под конкретные режимы.
\item Единый механизм восстановления контроллера после обнаружения через инфляцию ковариационной матрицы \RLS{} и временный сброс коэффициента забывания с линейным возвратом к номинальному значению.
\item Непрерывно модулируемая доверительная область между средним и внутренним слоями: радиус $\delta(\sigma_t)$ интерполируется по насыщенному сигналу тяжести $\sigma_t=\mathrm{clip}(s_t,0,1)$, где диагностическая статистика $s_t\in[0,10]$, что обеспечивает гладкое переключение между жёстким и мягким связыванием уровней.
\item Корректировка формул адаптации Сейджа-Хузы $Q,R$ с проекцией собственных значений на относительный нижний предел~--- устранение известного недостатка классической схемы Лу 2015~\cite{lu2015}, при котором матрицы могут терять положительную определённость в режиме разогрева.
\end{enumerate}

Новизна по существу состоит в \emph{архитектурном решении}: обнаружитель открытого мира и иерархическая инкрементальная адаптация заменяют классическую связку MMAE-банк + переключение контроллеров. Каждый отдельный компонент имеет литературные предтечи (\AINDI{}~\cite{aaindi2024}, \IADP{}~\cite{konatala2024}, фильтр Калмана с адаптацией Сейджа-Хузы~\cite{lu2015}, тест разладки \CUSUM{}~\cite{page1954}); вкладом работы является их компоновка в единый каскад с непрерывным сигналом тяжести и снятие допущения о фиксированном каталоге отказов.

Структура статьи: в разделе~\ref{sec:arch} описана архитектура \UFTC{} в полной шестислойной редакции и текущая трёхслойная рабочая реализация алгоритма. Раздел~\ref{sec:components} содержит детали трёх ключевых компонентов: $L_2$~--- \AINDI{}-обёртка с супер-твист-наблюдателем и ограниченной доверительной областью; $L_3$~--- \IADP{}-обёртка с инициируемым обнаружителем сбросом \RLS{}; параллельный обнаружитель \FDD{} на основе фильтра Калмана и теста \CUSUM{}. Раздел~\ref{sec:theory} приводит теоретические замечания: корректировка формул Сейджа-Хузы и качественное описание условий устойчивости композиции. В разделе~\ref{sec:exp} представлен численный эксперимент на нелинейной 6-DoF модели Boeing~747-100 со сценарием отказа крайнего левого двигателя. Сравнение проводится с двумя референсными подходами~--- разомкнутым контуром на статической балансировке и единым онлайн-актором событийно-триггерного DHP~\cite{mazaev2026etdhp}, реализованным в той же программной среде. Раздел~\ref{sec:disc} обсуждает применимость и ограничения \UFTC{} в текущей версии; раздел~\ref{sec:concl} формулирует выводы и направления продолжения.

% ====================================================================
\section{Архитектура UFTC}
\label{sec:arch}
% ====================================================================

\subsection{Полная шестислойная редакция}

В полной редакции \UFTC{} спроектирован как шестислойная композиция (рис.~\ref{fig:arch-full}):

\begin{itemize}
\item $L_4$~--- внешний контур на основе дистрибутивного актора-критика мягкого Soft Actor-Critic с фильтром риска CVaR; обрабатывает командные сигналы пилота либо системы планирования миссии и формирует эталон управления для среднего слоя~\cite{ma2020dsac}.
\item $L_3$~--- средний контур на инкрементальном ADP с TD($\lambda$): онлайн-идентификация локальной линейной модели $\Delta x_{t+1} = \tilde F_t \Delta x_t + \tilde G_t \Delta u_t$ через \RLS{}, оценка функции ценности $\tilde V_\pi$ и улучшение политики методом замкнутого решения квадратичной задачи~\cite{konatala2024,zhou2018}.
\item $L_2$~--- внутренний контур на активно-адаптивном INDI с супер-твист-наблюдателем неподмоделированного возмущения и переключателем режима «по угловой скорости / по углу атаки»~\cite{aaindi2024,sieberling2010}.
\item $L_1$~--- защитный фильтр на основе достижимости Гамильтона-Якоби с предобученной функцией ценности безопасной области полёта и сертифицированными липшицевыми границами через CROWN; обеспечивает невыход из envelope-границ на каждом такте~\cite{lygeros1999,bansal2017}.
\item Параллельный обнаружитель отказов \FDD{}~--- единственный номинальный фильтр Калмана с тестом разладки CUSUM по махаланобисовой дистанции инноваций; формирует бинарный сигнал \texttt{fault\_present} и непрерывный диагностический сигнал тяжести $s_t$.
\item Композитный векторный монитор Ляпунова~--- 5-компонентная функция $V_{\text{total}} = \sum_i w_i V_i$, проверяющая выполнение условий устойчивости каскада на каждом такте~\cite{khalil2002}.
\end{itemize}

\begin{figure}[H]
\centering
\resizebox{\linewidth}{!}{%
\begin{tikzpicture}[>=Stealth,
    layer/.style={rectangle,draw,rounded corners=2pt,minimum width=58mm,minimum height=10mm,align=center,font=\small},
    side/.style={rectangle,draw,rounded corners=2pt,minimum width=42mm,minimum height=12mm,align=center,font=\small,fill=red!8},
    monitor/.style={rectangle,draw,rounded corners=2pt,minimum width=42mm,minimum height=12mm,align=center,font=\small,fill=blue!8},
    sm/.style={font=\scriptsize,fill=white,inner sep=1pt}]
\node[layer] (L4) at (0,0) {$L_4$\\DSAC + CVaR};
\node[layer] (L3) at (0,-1.55) {$L_3$\\\IADP{} + RLS};
\node[layer] (L2) at (0,-3.10) {$L_2$\\\AINDI{} + super-twist};
\node[layer] (L1) at (0,-4.65) {$L_1$\\HJ-shield};
\node[layer] (act) at (0,-6.20) {Актуаторы / ОУ};
\node[side] (fdd) at (8.2,-2.15) {\FDD{}\\Kalman + CUSUM};
\node[monitor] (mon) at (8.2,-5.40) {Vector\\Lyapunov monitor};

\draw[->] (L4) -- node[sm,right] {$u_{\text{SAC}},\,\beta_t$} (L3);
\draw[->] (L3) -- node[sm,right] {$u_{\text{iADP}},\,\omega_{\text{ref}}$} (L2);
\draw[->] (L2) -- node[sm,right] {$u_{\text{INDI}}$} (L1);
\draw[->] (L1) -- node[sm,right] {$u_t$} (act);

\coordinate (diagbus) at (4.45,-2.15);
\draw[->] (fdd.west) -- node[sm,above] {$s_t,\sigma_t,p_f$} (diagbus);
\draw[densely dashed] (4.45,0) -- (4.45,-4.65);
\draw[->] (4.45,0) -- node[sm,above] {$p_f$} (L4.east);
\draw[->] (4.45,-1.55) -- node[sm,above] {$s_t$} (L3.east);
\draw[->] (4.45,-3.10) -- node[sm,above] {$\sigma_t$} (L2.east);
\draw[->] (4.45,-4.65) -- node[sm,above] {$s_t$} (L1.east);
\draw[->] (act.east) -- ++(8.10,0) |- node[sm,pos=0.27,right] {$x_t,u_{t-1}$} (fdd.east);
\draw[->,dashed] (act.east) -- ++(0.95,0) |- node[sm,pos=0.83,above] {$V_i$} (mon.west);
\end{tikzpicture}
}
\caption{Полная шестислойная архитектура \UFTC{}. Стрелки слева~--- иерархия команд $L_4 \to L_3 \to L_2 \to L_1 \to$ актуаторы. Параллельные блоки \FDD{} и vector Lyapunov monitor связаны со всеми слоями.}
\label{fig:arch-full}
\end{figure}

\subsection{Трёхслойная рабочая реализация}
\label{sec:implemented-scope}

В настоящей работе реализована рабочая трёхслойная редакция алгоритма, включающая три центральных компонента: внутренний $L_2$, средний $L_3$ и обнаружитель \FDD{}. Слой $L_4$ в экспериментальном контуре представлен единичным сигналом риска ($\beta_t \equiv 1$, выход \IADP{} проходит без модуляции), а $L_1$ HJ-shield и vector Lyapunov monitor рассматриваются как расширения полной архитектуры. Такое построение фиксирует уже работоспособный каскад на наиболее критичных интерфейсах (\FDD{} $\leftrightarrow L_3$ $\leftrightarrow L_2$) и задаёт алгоритм как единую дискретную САУ с завершённым вычислительным циклом.

Поверхность интерфейсов текущей реализации показана на рис.~\ref{fig:arch-implemented}. С точки зрения теории автоматического управления получившаяся структура~--- двухконтурная САУ с подчинённым регулированием: внешний контур $L_3$ формирует эталонные значения угловых скоростей по ошибке стабилизации режима полёта, внутренний контур $L_2$ отрабатывает эти эталоны на исполнительных механизмах. Параллельный канал \FDD{} реализует функцию надзорного автомата: следит за невязкой между моделью объекта и измерениями, и при выходе невязки за допустимые границы изменяет настройки регулирующих контуров (расширяет доверительную область в $L_2$, инициирует переидентификацию в $L_3$).

\begin{figure}[H]
\centering
\resizebox{\linewidth}{!}{%
\begin{tikzpicture}[>=Stealth,
    layer/.style={rectangle,draw,rounded corners=2pt,minimum width=64mm,minimum height=12mm,align=center,font=\small},
    fdd/.style={rectangle,draw,rounded corners=2pt,minimum width=44mm,minimum height=14mm,align=center,font=\small,fill=red!8},
    sm/.style={font=\scriptsize,fill=white,inner sep=1pt}]
\node[layer] (L3) at (0,0) {$L_3$ \IADP{}\\$x_t,r_t \mapsto (u_{\text{iADP}},\omega_{\text{ref}})$};
\node[layer] (L2) at (0,-1.80) {$L_2$ \AINDI{}\\$(\omega_{\text{ref}},\omega_t,u_{\text{iADP}},\sigma_t) \mapsto u_{\text{raw}}$};
\node[layer] (act) at (0,-3.60) {Актуаторы / летательный аппарат};
\node[fdd] (fdd) at (6.7,-1.80) {\FDD{}\\Kalman + CUSUM\\$(s_t,\text{fault})$};

\draw[->] (L3) -- node[sm,right] {$\omega_{\text{ref}},u_{\text{iADP}}$} (L2);
\draw[->] (L2) -- node[sm,right] {$u_t$} (act);
\draw[->] (act.east) -- ++(0.8,0) |- node[sm,pos=0.25,right] {$x_t,u_{t-1}$} (fdd.south);
\draw[->] (fdd.west) -- node[sm,above] {$s_t,\sigma_t$} (L2.east);
\draw[->] (fdd.north) |- node[sm,pos=0.25,above] {$\text{fault}$} (L3.east);
\end{tikzpicture}
}
\caption{Рабочая трёхслойная реализация \UFTC{}: слои $L_2$, $L_3$ и параллельный обнаружитель \FDD{}. Сигнал $s_t$ управляет сбросом \RLS{} в $L_3$, а $\sigma_t$ задаёт радиус доверительной области в $L_2$.}
\label{fig:arch-implemented}
\end{figure}

\subsection{Алгоритм работы UFTC за один такт дискретизации}
\label{sec:algorithm}

Для удобства восприятия приведём порядок вычислений, выполняемых контроллером \UFTC{} в течение одного такта дискретизации $\Delta t$. Все блоки работают синхронно от единого тактового генератора; задержка распространения сигнала между блоками не превышает одного такта (рис.~\ref{fig:timeline}).

\begin{figure}[H]
\centering
\resizebox{\linewidth}{!}{%
\begin{tikzpicture}[>=Stealth,
    tickstep/.style={rectangle,draw,rounded corners=1.5pt,minimum width=32mm,minimum height=11mm,align=center,font=\footnotesize},
    fdd/.style={tickstep,fill=red!8},
    l3/.style={tickstep,fill=blue!6},
    l2/.style={tickstep,fill=green!8},
    plant/.style={tickstep,fill=gray!12},
    sm/.style={font=\scriptsize,fill=white,inner sep=1pt}]
\node[plant] (s1) at (0,0) {1. Измерения\\$x_t,\omega_t$};
\node[fdd] (s2) at (4.0,0) {2. \FDD{}\\Kalman + CUSUM};
\node[l3] (s3) at (8.0,0) {3. $L_3$\\RLS + политика};
\node[l2] (s4) at (4.0,-1.90) {4. $L_2$\\\AINDI{} + наблюдатель};
\node[tickstep] (s5) at (8.0,-1.90) {5. Проекция\\$\mathrm{proj}_{\delta(\sigma_t)}$};
\node[plant] (s6) at (12.0,-1.90) {6. Подача\\$u_t$ на ОУ};

\draw[->] (s1) -- node[sm,above] {$x_t,u_{t-1}$} (s2);
\draw[->] (s2) -- node[sm,above] {$s_t,\sigma_t,\text{fault}$} (s3);
\draw[->] (s3.south) -- ++(0,-0.55) -- node[sm,above] {$\omega_{\text{ref}},u_{\text{iADP}}$} (4.0,-1.05) -- (s4.north);
\draw[->] (s4) -- node[sm,above] {$u_{\text{INDI,raw}}$} (s5);
\draw[->] (s5) -- node[sm,above] {$u_t$} (s6);
\draw[->,dashed] (s6.south) |- ++(-12.0,-0.85) -| node[sm,pos=0.28,below] {регистры такта $t+1$} (s1.south);
\end{tikzpicture}
}
\caption{Последовательность операций \UFTC{} в одном такте: 1)~получение измерений; 2)~обновление обнаружителя \FDD{}, формирование сигнала тяжести $s_t$ и насыщенной величины $\sigma_t$; 3)~предсказание управления \IADP{}, идентификация локальной модели через \RLS{}; 4)~предсказание управления \AINDI{} с компенсацией возмущений; 5)~ограничение по доверительной области с радиусом, зависящим от $\sigma_t$; 6)~подача результата на ОУ. Пунктирная линия~--- передача состояния на следующий такт.}
\label{fig:timeline}
\end{figure}

Развёрнутый псевдокод одного такта приведён в алгоритме~\ref{alg:uftc-step}.

\begin{center}
\begin{minipage}{0.96\linewidth}
\refstepcounter{algorithm}\label{alg:uftc-step}
\rule{\linewidth}{0.6pt}\\[2pt]
\textbf{Алгоритм~\thealgorithm.} Один такт работы \UFTC{} в трёхслойной реализации.\\[2pt]
\rule{\linewidth}{0.4pt}\\[3pt]
\textit{Вход:} измерение $x_t$ полного вектора состояния ОУ, угловые скорости $\omega_t = (p,q,r)^T$, эталон $r_t$ режима полёта, предыдущее управление $u_{t-1}$, накопленная статистика \CUSUM{} $S_{t-1}$, ковариация \RLS{} $\Phi_{t-1}$, оценка матриц $\tilde F_{t-1}, \tilde G_{t-1}$, фаза восстановления $\gamma_\text{RLS}(t-1)$.\\
\textit{Выход:} управление $u_t \in \mathbb{R}^{n_u}$, обновлённое состояние внутренних регистров.\\[5pt]
\textbf{1. Обнаружение и формирование сигнала тяжести (\FDD{}):}
\begin{enumerate}
\item[1.1.] Прогнозирование номинальной модели:
$\hat x^- \leftarrow (I+F)\hat x_{t-1} + Gu_{t-1}$, $P^- \leftarrow (I+F)P(I+F)^\top + Q$.
\item[1.2.] Вычисление инноваций: $\nu_t \leftarrow x_t - \hat x^-$, ковариация инноваций $S_\nu \leftarrow P^- + R$.
\item[1.3.] Махаланобисова дистанция: $d_t \leftarrow \nu_t^\top S_\nu^{-1} \nu_t$.
\item[1.4.] Тест разладки: $S_t \leftarrow \max(0,\, S_{t-1} + d_t - \mathrm{drift})$.
\item[1.5.] Сигнал тяжести: $s_t \leftarrow \mathrm{clip}(S_t/h_\text{alarm},\,0,\,10)$, $\sigma_t \leftarrow \mathrm{clip}(s_t,\,0,\,1)$, флаг $\text{fault}_t$ по гистерезисной логике.
\item[1.6.] Коррекция оценки: $K \leftarrow P^- S_\nu^{-1}$, $\hat x_t \leftarrow \hat x^- + K\nu_t$, $P \leftarrow (I-K)P^-(I-K)^\top + KRK^\top$.
\end{enumerate}
\textbf{2. Средний контур $L_3$ (\IADP{}):}
\begin{enumerate}
\item[2.1.] Если на текущем такте обнаружен передний фронт $\text{fault}_t$~--- инициировать сброс: $\Phi_t \leftarrow \Phi_t + \kappa I$, $\gamma_\text{RLS}(t) \leftarrow \gamma_\text{drop}$.
\item[2.2.] Иначе~--- линейное возвращение коэффициента забывания к номинальному значению.
\item[2.3.] Обновление \RLS{}-оценки $\tilde \Theta_t = [\tilde F_t;\,\tilde G_t]^\top$ по новой паре $(x_t,u_{t-1})$.
\item[2.4.] Решение задачи квадратичной оптимизации~\eqref{eq:iadp-policy}, формирование $u_\text{iADP}$.
\item[2.5.] Извлечение эталона угловых скоростей $\omega_\text{ref}$ по~\eqref{eq:omega-ref}.
\end{enumerate}
\textbf{3. Внутренний контур $L_2$ (\AINDI{}):}
\begin{enumerate}
\item[3.1.] Обновление супер-твист-наблюдателя~\eqref{eq:supertwist-eta}, \eqref{eq:supertwist-z}; коррекция измерения $\omega_\text{corr} \leftarrow \omega_t - \hat\delta\Delta t$.
\item[3.2.] Расчёт сырого управления $u_\text{INDI,raw}$ методом \INDI{} по эталону $\omega_\text{ref}$ и скорректированному измерению $\omega_\text{corr}$.
\end{enumerate}
\textbf{4. Доверительная область:}
\begin{enumerate}
\item[4.1.] Радиус $\delta(\sigma_t)$ по~\eqref{eq:trust-radius}.
\item[4.2.] Проекция: $u_t \leftarrow u_\text{iADP} + \mathrm{clip}_\delta(u_\text{INDI,raw} - u_\text{iADP})$ по~\eqref{eq:trust-region}.
\end{enumerate}
\textbf{5. Подача управления на ОУ:} $u_t \to$ актуаторы; обновление шага $t \leftarrow t+1$.\\[2pt]
\rule{\linewidth}{0.6pt}
\end{minipage}
\end{center}

Структура алгоритма соответствует классической схеме каскадного регулятора с подчинённым регулированием по угловой скорости (внутренний контур $L_2$) и по угловому положению / режиму полёта (внешний контур $L_3$). Параллельный канал \FDD{} формирует надзорный (диагностический) сигнал, не разрывая основную цепь регулирования и не внося задержку на исполнительные механизмы. При срабатывании обнаружителя диагностическая статистика $s_t$ и её насыщенная форма $\sigma_t$ модулируют параметры регуляторов, не переключая их структуру~--- это сохраняет непрерывность управления и устраняет ударные процессы при бамплесс-переходе.

\paragraph{Инженерная трактовка сигнального тракта.} С точки зрения разработчика САУ \UFTC{} представляет собой супервизорно-модулируемый каскадный регулятор. Основной управляющий канал имеет последовательную структуру
\[
x_t,\ r_t \;\longrightarrow\; L_3 \;\longrightarrow\; L_2 \;\longrightarrow\; \mathrm{proj}_{\delta(\sigma_t)} \;\longrightarrow\; u_t \;\longrightarrow\; \text{ОУ},
\]
а канал \FDD{} включён параллельно как диагностический супервизор. Он не стоит последовательно в цепи управления и поэтому не может сам по себе разорвать контур стабилизации. Его выходы используются только как параметры настройки: $s_t$ задаёт уровень диагностической аномалии, $\sigma_t$ ограничивает этот уровень диапазоном, пригодным для модуляции регуляторов, а $\text{fault}_t$ является дискретным событием для одноразового сброса идентификатора.

Алгоритм реализуется как дискретная САУ с регистровой логикой. В начале такта фиксируются измерения $x_t,\omega_t$ и регистры предыдущего шага: $u_{t-1}$, $\hat x_{t-1}$, $P_{t-1}$, $S_{t-1}$, $\Phi_{t-1}$, $\tilde F_{t-1}$, $\tilde G_{t-1}$, $\gamma_{\text{RLS}}(t-1)$. Все вычисления внутри такта используют этот снимок состояния. Команда $u_t$ становится входом ОУ только в конце такта, а в фильтр обнаружителя она попадёт как $u_{t-1}$ на следующем шаге. Такое соглашение устраняет алгебраическую петлю между фильтром, регулятором и объектом.

\paragraph{Назначение контуров.} Внутренний слой $L_2$ трактуется как быстрый демпфирующий контур по угловым скоростям. Его задача~--- удерживать локальную угловую динамику и парировать быстрые возмущения в каналах крена, тангажа и рыскания. Средний слой $L_3$ является более медленным адаптивно-идентификационным контуром: он уточняет локальную инкрементальную модель ОУ и формирует базовую команду $u_{\text{iADP}}$ вместе с эталоном $\omega_{\text{ref}}$ для внутреннего контура. Разделение ролей соответствует стандартному принципу подчинённого регулирования: быстрый внутренний контур обеспечивает исполнительную устойчивость, медленный внешний контур меняет рабочую точку и перераспределяет управление при изменении эффективной динамики.

Обнаружитель \FDD{} является надзорным каналом. На каждом шаге он сравнивает фактическое измерение с прогнозом номинальной модели, вычисляет инновацию $\nu_t$ и нормирует её через ковариацию инноваций. Махаланобисова дистанция $d_t$ выступает безразмерной мерой рассогласования «объект~--- номинальная модель». Накопитель \CUSUM{} не реагирует на одиночный шумовой выброс как на отказ, а интегрирует устойчивое смещение статистики. Поэтому выход $\text{fault}_t$ имеет гистерезис и cooldown, а непрерывный $s_t$ даёт регулятору ранний подпороговый признак ухудшения динамики.

\paragraph{Последовательность работы на такте.} Первым выполняется диагностический прогноз: фильтр Калмана строит $\hat x^{-}$ из предыдущей оценки и предыдущего управления $u_{t-1}$. Это принципиально важно для реализации в бортовом ПО: обнаружитель работает на уже применённой команде и не создаёт прямой зависимости от ещё не рассчитанного $u_t$. После коррекции фильтра получаются $d_t$, $S_t$, $s_t$, $\sigma_t$ и дискретное состояние автомата обнаружения.

Далее отрабатывает $L_3$. Если на текущем такте зарегистрирован нарастающий фронт $\text{fault}_t$, то выполняется не переключение регулятора, а изменение режима идентификатора: ковариационная матрица \RLS{} раздувается, а коэффициент забывания временно уменьшается. Для разработчика САУ это аналог перевода оценивателя в режим высокой податливости к новым данным после структурного изменения объекта. При отсутствии фронта отказа \RLS{} продолжает штатную идентификацию, а $\gamma_{\text{RLS}}$ плавно возвращается к номинальному значению. На выходе $L_3$ получается базовое абсолютное управление $u_{\text{iADP}}$ и командный сигнал для быстрого контура $\omega_{\text{ref}}$.

После этого работает $L_2$. Супер-твист-наблюдатель оценивает быструю неподмоделированную составляющую в канале угловых скоростей и формирует поправку $\hat\delta$. \AINDI{} использует скорректированное измерение $\omega_{\text{corr}}$ и эталон $\omega_{\text{ref}}$, чтобы получить сырой управляющий вектор $u_{\text{INDI,raw}}$. Этот вектор может быть более агрессивным, чем команда среднего слоя, поэтому он не подаётся на ОУ напрямую.

Финальным элементом является доверительная область вокруг $u_{\text{iADP}}$. Проекция $\mathrm{proj}_{\delta(\sigma_t)}$ ограничивает расхождение между быстрым INDI-действием и базовой командой среднего контура:
\[
\|u_t-u_{\text{iADP}}\| \leq \delta(\sigma_t).
\]
В номинальном режиме $\sigma_t\approx 0$, поэтому внутренний контур жёстко привязан к среднему слою и не вносит лишних управляющих движений. При росте диагностической аномалии радиус увеличивается, и $L_2$ получает больший оперативный ресурс для парирования возмущения. При этом структура закона управления остаётся той же, поэтому переход является бамплесс: нет скачка из-за выбора другого регулятора или другого банка наблюдателей.

\paragraph{Режимы работы.} В штатном режиме $S_t<h_{\text{clear}}$, $\sigma_t$ близка к нулю, \RLS{} работает с номинальным забыванием, а доверительная область имеет минимальный радиус. В подпороговом режиме $h_{\text{clear}}\leq S_t<h_{\text{alarm}}$ бинарный отказ ещё не объявлен, но $\sigma_t$ уже расширяет допустимое отклонение $L_2$ от $u_{\text{iADP}}$. В режиме тревоги $S_t\geq h_{\text{alarm}}$ на первом фронте выполняется сброс идентификатора \RLS{}; последующие такты тревоги не должны повторно раздувать $\Phi$, чтобы не получить самовозбуждение оценивателя. В режиме восстановления статистика \CUSUM{} спадает ниже $h_{\text{clear}}$, но выход автомата удерживается cooldown-интервалом, что предотвращает дребезг дискретного сигнала при работе около порога.

Таким образом, алгоритм не является набором независимых адаптивных блоков. Это единая дискретная САУ с основным каскадом регулирования, параллельной диагностикой, одноразовыми событийными воздействиями на идентификатор и непрерывной параметрической модуляцией быстрого контура.

% ====================================================================
\section{Компоненты рабочей реализации}
\label{sec:components}
% ====================================================================

\subsection{\texorpdfstring{Средний слой $L_3$: IADPMiddle}{Средний слой L3: IADPMiddle}}

В основе среднего слоя~--- агент \IADP{}, реализующий схему Konatala et al.~\cite{konatala2024}: онлайн-идентификация инкрементальной локальной линейной модели через \RLS{} с постоянным коэффициентом забывания, оценка ядра функции ценности $\tilde V_\pi(X) = X^\top \tilde P X$ методом наименьших квадратов на скользящем окне переходов и закрытое решение задачи минимизации стоимости-к-цели:
\begin{equation}
\Delta\delta_t = -(R + \gamma \tilde G_t^\top \tilde P \tilde G_t)^{-1}\bigl[R\,\delta_{t-1} + \gamma \tilde G_t^\top \tilde P \hat X_t + \gamma \tilde G_t^\top \tilde P \tilde F_t \Delta \hat X_t\bigr].
\label{eq:iadp-policy}
\end{equation}
где $X_t = [x_t;\,x_t^r]$ — расширенное состояние, объединяющее наблюдение объекта и эталонный сигнал, $R$ и $Q$ — весовые матрицы стоимости-к-шагу, $\gamma$ — коэффициент дисконтирования.

Обёртка \texttt{IADPMiddle} расширяет базовый агент двумя механизмами:

\paragraph{Механизм инициируемого обнаружителем сброса \RLS{}.} На каждом такте обёртка получает выход \FDD{} в виде кортежа $(\text{fault\_present},\,s_t)$, где $s_t\in[0,10]$ — нормированная статистика \CUSUM{}, отражающая тяжесть отказа. На переднем фронте сигнала $\text{fault\_present}$ (нарастающий перепад) выполняется триггер сброса:
\begin{align}
\Phi_t  &\leftarrow \Phi_t + \kappa I, \quad \kappa = 100, \label{eq:phi-inflate}\\
\gamma_{\text{RLS}}(t) &\leftarrow \gamma_{\text{drop}} = 0{,}9. \label{eq:gamma-drop}
\end{align}
Раздувание ковариации $\Phi$ (\ref{eq:phi-inflate}) возвращает \RLS{} в режим высокой неопределённости, в котором последующие обновления способны быстро двигать оценку параметров $\tilde\Theta = [\tilde F;\tilde G]^\top$ к новым значениям, соответствующим повреждённой динамике. Снижение коэффициента забывания $\gamma_{\text{RLS}}$ с номинального $0{,}99$ до $0{,}9$ ускоряет «забывание» предотказной модели. После триггера запускается линейный возврат $\gamma_{\text{RLS}}$ к номинальному значению за $N_{\text{rec}} = 500$ шагов:
\begin{equation}
\gamma_{\text{RLS}}(t+k) = \gamma_{\text{drop}} + (\gamma_{\text{nom}} - \gamma_{\text{drop}}) \cdot \frac{k}{N_{\text{rec}}}, \quad k\in[0,N_{\text{rec}}].
\label{eq:gamma-recover}
\end{equation}

\paragraph{Извлечение эталона угловых скоростей.} \IADP{}-выход~--- абсолютное управление $u_{\text{iADP}}$ в исходных физических единицах. Внутренний слой $L_2$, в свою очередь, ожидает эталон угловой скорости $\omega_{\text{ref}}$. Обёртка вычисляет эталон угловой скорости методом упреждения по текущей ошибке состояния:
\begin{equation}
\omega_{\text{ref}}(t) = \frac{r_t^{\omega} - \omega_{\text{meas}}(t)}{\Delta t_{\text{LH}}},
\label{eq:omega-ref}
\end{equation}
где $r_t^{\omega}$~--- компоненты эталонного сигнала, отвечающие угловым скоростям, $\Delta t_{\text{LH}}$~--- горизонт упреждения. При отсутствии явного указания эталон $\omega_{\text{ref}}$ интерпретируется как первичный (нескорректированный) сигнал.

\subsection{\texorpdfstring{Внутренний слой $L_2$: WrappedAAINDI}{Внутренний слой L2: WrappedAAINDI}}

Внутренний слой реализован как обёртка над агентом \AINDI{}~\cite{aaindi2024}, дополненная тремя механизмами:

\paragraph{Супер-твист-наблюдатель неподмоделированного возмущения.} По схеме высокопорядкового скользящего режима Леванта~\cite{levant1993}:
\begin{align}
\dot \eta &= -k_1 |\eta|^{1/2} \mathrm{sgn}(\eta) + z, \label{eq:supertwist-eta}\\
\dot z &= -k_2 \mathrm{sgn}(\eta),  \label{eq:supertwist-z}
\end{align}
где $\eta$~--- скользящая переменная, $z$~--- интегральная компонента, $k_1=3{,}0$, $k_2=1{,}5$. Выход $\hat\delta = \eta$ инжектируется в измерение угловой скорости через коррекцию $\omega_{\text{corr}} = \omega_{\text{meas}} - \hat\delta\,\Delta t$, что компенсирует высокочастотные неподмоделированные эффекты (несимметричный момент рыскания, индуцированный отказом двигателя, аэроупругие колебания).

\paragraph{Гистерезисный переключатель режима «угловая скорость / угол атаки».} Условие переключения с rate-INDI на angle-INDI задано порогом $\alpha_{\text{thr}}=25\degr$ с гистерезисной зоной $\Delta\alpha=5\degr$:
\begin{equation}
\text{mode}(t) = \begin{cases}
\text{angle}, & \alpha(t) > \alpha_{\text{thr}} \text{ и mode}(t-1)=\text{rate}, \\
\text{rate}, & \alpha(t) < \alpha_{\text{thr}} - \Delta\alpha \text{ и mode}(t-1)=\text{angle}, \\
\text{mode}(t-1), & \text{иначе}.
\end{cases}
\label{eq:mode-switch}
\end{equation}
В текущей реализации выбор режима определяет диагностическую метку, а закон управления остаётся единым (rate-INDI); полноценное подключение angle-INDI рассматривается как дальнейшее расширение внутреннего слоя.

\paragraph{Ограниченная доверительная область с непрерывной модуляцией.} Финальное управление $u_{\text{INDI}}$ ограничивается шаром радиуса $\delta(\sigma_t)$ вокруг абсолютного управления $u_{\text{iADP}}$, поступающего со среднего слоя:
\begin{equation}
u_{\text{INDI}} = u_{\text{iADP}} + \min\!\left(\delta(\sigma_t),\; \|u_{\text{INDI,raw}} - u_{\text{iADP}}\|\right) \cdot \frac{u_{\text{INDI,raw}} - u_{\text{iADP}}}{\|u_{\text{INDI,raw}} - u_{\text{iADP}}\|},
\label{eq:trust-region}
\end{equation}
где $u_{\text{INDI,raw}}$~--- сырой выход \AINDI{}-агента после супер-твист-коррекции, а радиус интерполируется по насыщенному сигналу тяжести $\sigma_t=\mathrm{clip}(s_t,0,1)$:
\begin{equation}
\delta(\sigma_t) = \delta_{\text{nom}} + (\delta_{\text{fault}} - \delta_{\text{nom}}) \cdot \sigma_t.
\label{eq:trust-radius}
\end{equation}
Под номинальным режимом ($\sigma_t=0$) радиус $\delta_{\text{nom}}$ невелик ($0{,}05$ в нормированных единицах), и $L_2$ жёстко связан с решением \IADP{}-агента; при достижении порога тревоги ($\sigma_t=1$) радиус расширяется до $\delta_{\text{fault}}=0{,}40$, давая среднему слою свободу для переадаптации. Гладкая интерполяция вместо дискретного переключения исключает разрывы команды управления при переходе через границу \texttt{fault\_present}.

\subsection{Параллельный обнаружитель FDD}

Обнаружитель \FDD{} (рис.~\ref{fig:fdd}) состоит из единственного номинального фильтра Калмана \texttt{NominalKalman} и теста разладки \CUSUM{} с гистерезисом \texttt{ChangePointDetector}. Фильтр настроен на номинальную динамику и работает по инкрементальной линеаризованной модели
\begin{equation}
x_{t+1} \approx (I + F)\,x_t + G\,u_t,
\label{eq:kalman-model}
\end{equation}
где $F$, $G$~--- начальные оценки локальных системной и управляющей матриц (начальная инициализация), которые могут быть либо предоставлены пользователем (если линеаризация в окрестности номинального режима известна), либо извлечены из \RLS{}-оценки \IADP{} по окончании заданного интервала разогрева.

\begin{figure}[H]
\centering
\resizebox{\linewidth}{!}{%
\begin{tikzpicture}[>=Stealth,
    blk/.style={rectangle,draw,rounded corners=2pt,minimum width=36mm,minimum height=13mm,align=center,font=\small},
    io/.style={rectangle,draw,rounded corners=2pt,minimum width=28mm,minimum height=9mm,align=center,font=\small,fill=gray!10},
    sm/.style={font=\scriptsize,fill=white,inner sep=1pt}]
\node[io] (inp) at (0,1.95) {Измерения\\$x_t,u_{t-1}$};
\node[blk] (kf) at (0,0) {NominalKalman\\Sage--Husa\\PSD projection};
\node[blk] (cpd) at (5.8,0) {ChangePointDetector\\CUSUM + hysteresis};
\node[blk] (out) at (11.2,0) {FDDOutput\\$s_t,\sigma_t,\text{fault}$};

\draw[->] (inp) -- (kf);
\draw[->] (kf) -- node[sm,above] {$\nu_t,\ S_\nu$} node[sm,below] {$d_t$} (cpd);
\draw[->] (cpd) -- node[sm,above] {$S_t$} (out);
\end{tikzpicture}
}
\caption{Структура обнаружителя \FDD{}: измерение $x_{\text{meas}}$ и предыдущее управление $u_{t-1}$ поступают в фильтр Калмана, нормированная махаланобисова дистанция инноваций $d_t = \nu_t^\top S_t^{-1} \nu_t$ подаётся на тест \CUSUM{}, выход~--- бинарный сигнал \texttt{fault\_present} и непрерывный диагностический сигнал тяжести $s_t$.}
\label{fig:fdd}
\end{figure}

\paragraph{Фильтр Калмана с адаптацией Сейджа-Хузы.} Для номинальной динамики $H = I$ (полнопараметрическое наблюдение состояния) и $F$, $G$ как инкрементальных матриц, шаги фильтра имеют стандартную форму:
\begin{align}
\hat x^- &= \hat x + F \hat x + G u_{t-1}, \\
P^- &= (I+F) P (I+F)^\top + Q, \\
\nu  &= x_{\text{meas}} - \hat x^-, \quad S = P^- + R, \\
K  &= P^- S^{-1}, \\
\hat x^+ &= \hat x^- + K\nu, \quad P^+ = (I-K) P^- (I-K)^\top + K R K^\top.
\end{align}

Адаптация ковариаций $Q$ и $R$ по схеме Сейджа--Хузы~\cite{sage1969,lu2015}:
\begin{align}
R^+_{\text{Lu}} &= (1-\alpha_R)\,R + \alpha_R \bigl(\nu\nu^\top - P^-\bigr), \label{eq:lu-R}\\
Q^+_{\text{Lu}} &= (1-\alpha_Q)\,Q + \alpha_Q \Bigl(K\,(\nu\nu^\top)\,K^\top + P^+ - (I+F)\,P\,(I+F)^\top\Bigr). \label{eq:lu-Q}
\end{align}

\paragraph{Тест разладки CUSUM с гистерезисом.} На каждом такте по выходу фильтра Калмана вычисляется махаланобисова дистанция $d_t$. Под нулевой гипотезой $H_0$ номинальной динамики выполняется $d_t \sim \chi^2_n$ с математическим ожиданием $n$. Тест \CUSUM{} обновляет накопленную статистику:
\begin{equation}
S_t = \max(0, S_{t-1} + d_t - \mathrm{drift}),
\label{eq:cusum}
\end{equation}
где $\mathrm{drift} = 1{,}5\,n$ — параметр сноса, обеспечивающий отрицательное математическое приращение под $H_0$ (в отличие от классического $\mathrm{drift} = n$, при котором тест представляет собой мартингейл и неизбежно пересекает любой конечный порог за конечное время). Бинарный выход:
\begin{equation}
\text{fault\_present}(t) = \begin{cases}
\text{True}, & S_t > h_{\text{alarm}} \text{ и } t > t_{\text{cooldown}}, \\
\text{False}, & S_t < h_{\text{clear}} \text{ и } t > t_{\text{cooldown}}, \\
\text{fault\_present}(t-1), & \text{иначе}.
\end{cases}
\label{eq:cusum-decision}
\end{equation}
Гистерезис $h_{\text{alarm}} > h_{\text{clear}}$ совместно с минимальным интервалом охлаждения $t_{\text{cooldown}}$ исключает чаттеринг сигнала на границе порога.

\paragraph{Сигнал тяжести.} Непрерывный скаляр тяжести $s_t$ и его насыщенная форма $\sigma_t$, передаваемые на $L_2$ и $L_3$:
\begin{equation}
s_t = \mathrm{clip}\!\left(\frac{S_t}{h_{\text{alarm}}},\, 0,\, 10\right).
\label{eq:severity}
\end{equation}
\begin{equation}
\sigma_t = \mathrm{clip}\!\left(s_t,\,0,\,1\right).
\label{eq:severity-saturated}
\end{equation}
$s_t = 1$ соответствует моменту пересечения порога; $s_t > 1$~--- продолжающемуся росту накопленной аномалии. Величина $\sigma_t$ используется там, где параметр регулятора должен оставаться ограниченным физическим диапазоном. Использование непрерывного сигнала вместо бинарного флага позволяет $L_2$ и $L_3$ реагировать пропорционально серьёзности отклонения.

\paragraph{Состояния обнаружителя как конечного автомата.} В терминах теории автоматов обнаружитель \FDD{} реализует автомат с четырьмя состояниями (рис.~\ref{fig:fdd-states}): \textsc{Норма} (накопленная статистика ниже порога очистки), \textsc{Подозрение} (накопление, но порог тревоги ещё не пересечён), \textsc{Тревога} (фактически зафиксированный отказ), \textsc{Восстановление} (период охлаждения после спада статистики). Переходы между состояниями определяются положением накопителя $S_t$ относительно порогов $h_\text{alarm}$ и $h_\text{clear}$ и счётчиком охлаждения $t_\text{cooldown}$. Гистерезис $h_\text{alarm} > h_\text{clear}$ исключает чаттеринг (быстрые переключения сигнала на границе порога) при шумовых выбросах статистики. Бинарный выход $\text{fault\_present}$ принимает значение \texttt{True} только в состоянии \textsc{Тревога}, что упрощает интерфейс с верхними слоями САУ.

\begin{figure}[H]
\centering
\resizebox{0.98\linewidth}{!}{%
\begin{tikzpicture}[>=Stealth,
    state/.style={rectangle,draw,rounded corners=2pt,minimum width=34mm,minimum height=14mm,align=center,font=\small},
    norm/.style={state,fill=green!8},
    susp/.style={state,fill=yellow!12},
    alrm/.style={state,fill=red!12},
    rec/.style={state,fill=blue!8},
    sm/.style={font=\scriptsize,fill=white,inner sep=1pt}]
\node[norm] (n) at (0,0) {\textsc{Норма}\\$S_t<h_\text{clear}$\\fault=False};
\node[susp] (s) at (4.5,0) {\textsc{Подозрение}\\$h_\text{clear}\le S_t<h_\text{alarm}$\\fault=False};
\node[alrm] (a) at (9.0,0) {\textsc{Тревога}\\$S_t\ge h_\text{alarm}$\\fault=True};
\node[rec] (r) at (4.5,-2.2) {\textsc{Восстановление}\\$S_t<h_\text{clear}$\\cooldown};

\draw[->] (n) -- node[sm,above] {$S_t$ растёт} (s);
\draw[->] (s) -- node[sm,above] {$S_t\ge h_\text{alarm}$} (a);
\draw[->] (s.south) -- node[sm,left] {$S_t<h_\text{clear}$} (r.north);
\draw[->] (a.south) |- node[sm,pos=0.30,right] {$S_t<h_\text{clear}$} (r.east);
\draw[->] (r.west) -| node[sm,pos=0.25,below] {$t\ge t_\text{cooldown}$} (n.south);
\end{tikzpicture}
}
\caption{Конечный автомат обнаружителя \FDD{}: четыре состояния и переходы между ними по уровням накопителя $S_t$ относительно порогов $h_\text{alarm}, h_\text{clear}$ и счётчика охлаждения. Гистерезис $h_\text{alarm} > h_\text{clear}$ устраняет чаттеринг при шумовых выбросах.}
\label{fig:fdd-states}
\end{figure}

% ====================================================================
\section{Теоретические замечания}
\label{sec:theory}
% ====================================================================

\subsection{Корректировка формул Сейджа-Хузы}

Классические формулы~\eqref{eq:lu-R} и~\eqref{eq:lu-Q} страдают известным недостатком: при разогреве (когда $P^-$ велика, инновации малы) выражение $\nu\nu^\top - P^-$ неотрицательно определено только при условии $\|\nu\|^2 \geq \mathrm{tr}(P^-)$, что в первые шаги фильтра не выполняется. Аналогично для $Q$: $P^+ - (I+F)P(I+F)^\top$ не имеет гарантии положительной определённости. Без специальной обработки матрицы $R$ и $Q$ могут терять симметрию и положительную определённость, что приводит к расходимости фильтра.

В настоящей работе предложена корректировка через проекцию собственных значений на относительный нижний предел:
\begin{equation}
M^{\text{PSD}} = V \cdot \mathrm{diag}(\max(\lambda_i,\,\lambda_{\text{floor}})) \cdot V^\top,
\label{eq:psd-proj}
\end{equation}
где $M = V \Lambda V^\top$~--- спектральное разложение симметризованной матрицы $\frac{1}{2}(M + M^\top)$, а нижний предел $\lambda_{\text{floor}}$ задан относительно начальных значений $Q_0$, $R_0$:
\begin{equation}
\lambda_{\text{floor}}^{R} = \max(10^{-9},\, 0{,}01 \cdot \lambda_{\min}(R_0)), \quad
\lambda_{\text{floor}}^{Q} = \max(10^{-9},\, 0{,}01 \cdot \lambda_{\min}(Q_0)).
\label{eq:psd-floor}
\end{equation}
Относительный нижний предел физически интерпретируется как минимальная неустранимая неопределённость, заданная конструктивной точностью датчика и квантованием процессной модели; абсолютный нижний предел $10^{-9}$ обеспечивает численную стабильность при произвольно сингулярных начальных значениях. Корректированные формулы:
\begin{align}
R^+ &= \mathrm{PSD\_proj}\!\bigl[(1-\alpha_R) R + \alpha_R (\nu\nu^\top - P^-)\bigr], \label{eq:R-corrected}\\
Q^+ &= \mathrm{PSD\_proj}\!\bigl[(1-\alpha_Q) Q + \alpha_Q (K\nu\nu^\top K^\top + P^+ - (I+F)P(I+F)^\top)\bigr]. \label{eq:Q-corrected}
\end{align}

\paragraph{Утверждение 1 (сохранение положительной определённости).} \emph{Если начальные ковариации $R_0, Q_0 \succ 0$, то для произвольной траектории фильтра, использующего обновления~\eqref{eq:R-corrected}, \eqref{eq:Q-corrected}, выполняется $R(t), Q(t) \succ 0$ для всех $t \geq 0$.}

\emph{Доказательство.} По построению оператора $\mathrm{PSD\_proj}$ из~\eqref{eq:psd-proj} все собственные значения результата ограничены снизу $\lambda_{\text{floor}} > 0$. Симметричность сохраняется в силу симметризации $\frac{1}{2}(M+M^\top)$ перед спектральным разложением. Следовательно $R(t), Q(t) \succeq \lambda_{\text{floor}} I \succ 0$. $\square$

\subsection{Условия устойчивости реализованного каскада}

Полное доказательство ограниченности замкнутой системы $L_2 + L_3 + \FDD{}$ по принципу Беллмана и матрицы Гурвица-Метцлера~\cite{khalil2002} требует выполнения следующих условий:
\begin{enumerate}
\item \textbf{$L_3$ \IADP{}.} Условия Konatala 2024~\cite{konatala2024} ограниченности (uniformly ultimately bounded, UUB) для квадратичной функции ценности и закрытого решения~\eqref{eq:iadp-policy}.
\item \textbf{$L_2$ \AINDI{}.} Условия Wang 2019~\cite{wang2019} устойчивости INDI с обобщением на неминимально-фазовые объекты через переключение rate-INDI/angle-INDI.
\item \emph{Время задержки переключения после обнаружения.} Минимальное время удержания режима после смены $\text{fault\_present}$ должно превышать $\tau_d > \log\mu / \lambda_{\min}$, где $\mu$ и $\lambda_{\min}$~--- параметры устойчивости каждого режима~\cite{liberzon2003}.
\item \emph{Сингулярное возмущение.} Соотношение скоростей адаптации внешнего и внутреннего слоёв $\epsilon = \alpha_{L_2} / \alpha_{L_3} \leq \epsilon^* \approx 0{,}1$.
\item \emph{Векторная функция Ляпунова.} Матрица $M$ с элементами $M_{ii} = -\alpha_i$, $M_{ij} = \beta_{ij}$ ($i \neq j$) гурвицева в смысле Метцлера.
\end{enumerate}

В текущей трёхслойной реализации условия 1 и 2 наследуются от исходных работ; условия 3-5 требуют дальнейшего формального доказательства, которое выходит за рамки настоящей статьи и составляет предмет отдельной теоретической публикации.

% ====================================================================
\section{Численный эксперимент}
\label{sec:exp}
% ====================================================================

\subsection{Постановка}

\paragraph{Нелинейный объект управления.} Эксперимент проводится на нелинейной 6-DoF модели Boeing~747-100 в крейсерском режиме~\cite{stevens2003,heffley1972}: эшелон $h_0 = 6096$~м (FL200), истинная воздушная скорость $V_0 \approx 205{,}4$~м/с (674~фут/с). Расчётная балансировочная конфигурация~--- симметричный полёт с тангажем $\vartheta_0 = +3{,}603\degr$, рулём высоты $\delta_{\text{в},0} = -0{,}722\degr$ и положением органа управления тягой $\delta_{\text{т},0} = 0{,}555$.

В терминах разработки САУ самолёт рассматривается как нелинейный многосвязный ОУ с двенадцатимерным состоянием
\begin{equation}
x =
\begin{bmatrix}
\boldsymbol v_b^T & \boldsymbol\omega_b^T & \boldsymbol\Theta^T & \boldsymbol r_e^T
\end{bmatrix}^{T},
\label{eq:b747-state}
\end{equation}
где $u_b,v_b,w_b$~--- проекции скорости центра масс на оси связанной системы координат; $p,q,r$~--- угловые скорости крена, тангажа и рыскания; $\varphi,\vartheta,\psi$~--- углы Эйлера; $x_e,y_e,z_e$~--- координаты в земной системе NED, причём высота $h=-z_e$. Вектор управляющих воздействий
\begin{equation}
u_c = \begin{bmatrix}\dv & \de & \dn & \dt\end{bmatrix}^{T}
\label{eq:b747-control}
\end{equation}
содержит отклонение руля высоты, элеронов, руля направления и нормированное положение РУД $\dt\in[0,1]$. Полная модель интегрирует все компоненты~\eqref{eq:b747-state}; в эксперименте \UFTC{} строит обратную связь по девяти регуляционным координатам~\eqref{eq:uftc-state-b747}, а внутренний \AINDI{}-контур выделяет из них подмножество $[p,q,r]$ для демпфирования угловых скоростей.

Правая часть ОУ записывается в стандартной форме пространственной динамики твёрдого тела:
\begin{align}
\dot{\boldsymbol v}_b &= \frac{1}{m}\bigl(\boldsymbol F_{\text{aero}}(x,u_c)+\boldsymbol F_{\text{eng}}(x,\dt,\mu)\bigr)
    + \boldsymbol g_b(\varphi,\vartheta) - \boldsymbol\omega_b\times \boldsymbol v_b, \\
\mathbf I\,\dot{\boldsymbol\omega}_b &=
    \boldsymbol M_{\text{aero}}(x,u_c)+\boldsymbol M_{\text{eng}}(x,\dt,\mu)
    - \boldsymbol\omega_b\times(\mathbf I\boldsymbol\omega_b), \\
\dot{\boldsymbol\Theta} &= \mathcal T(\varphi,\vartheta)\boldsymbol\omega_b,\qquad
\dot{\boldsymbol r}_e = C_b^e(\varphi,\vartheta,\psi)\boldsymbol v_b,
\label{eq:b747-6dof}
\end{align}
где $\boldsymbol v_b=[u_b,v_b,w_b]^T$, $\boldsymbol\omega_b=[p,q,r]^T$, $\boldsymbol\Theta=[\varphi,\vartheta,\psi]^T$, $\boldsymbol r_e=[x_e,y_e,z_e]^T$. Нелинейность ОУ обусловлена кинематикой Эйлера, гироскопическими членами $\boldsymbol\omega_b\times(\mathbf I\boldsymbol\omega_b)$, зависимостью динамического напора $\bar q=\frac{1}{2}\rho(h)V^2$ от скорости и высоты, а также аэродинамическими коэффициентами, вычисляемыми как локальное разложение по $\alpha$, $\beta$, $M$, угловым скоростям и отклонениям рулей:
\begin{equation}
C_\ast = C_{\ast 0}
+ C_{\ast \alpha}(\alpha-\alpha_0)
+ C_{\ast \beta}\beta
+ C_{\ast p}\frac{pb}{2V}
+ C_{\ast q}\frac{q\bar c}{2V}
+ C_{\ast r}\frac{rb}{2V}
+ C_{\ast \delta}u_c
+ C_{\ast M}(M-M_0).
\label{eq:b747-aero-build}
\end{equation}
Здесь $\alpha=\arctan(w_b/u_b)$, $\beta=\arcsin(v_b/V)$, $V=\sqrt{u_b^2+v_b^2+w_b^2}$; размерные силы и моменты получаются умножением безразмерных коэффициентов на $\bar q S$, $\bar q S b$ и $\bar q S\bar c$.

\paragraph{Моделирование отказа двигателя.} Подсистема повреждений не добавляет вектору состояния явный признак отказа и не переключает модель на заранее подготовленный линейный режим. Отказ задаётся через вектор эффективности четырёх двигателей
\begin{equation}
\mu(t)=
\begin{cases}
[1,\,1,\,1,\,1]^T, & t < t_f,\\
[0,\,1,\,1,\,1]^T, & t \ge t_f,
\end{cases}
\qquad t_f=10~\text{с}.
\label{eq:b747-engine-failure}
\end{equation}
Это соответствует пресету \texttt{LEFT\_OUTER\_ENGINE\_FAILURE}: двигатель~№\,1, крайний левый, полностью теряет тягу, а двигатели~№\,2--№\,4 остаются работоспособными.

Тяга каждого двигателя вычисляется из установленной модели Pratt \& Whitney JT9D-7:
\begin{equation}
T_i = \mu_i\,T_{\text{SLS}}^{(1)}
\bigl(1-0{,}49\sqrt{M}\bigr)\sigma(h)^{0{,}7}
\bigl(0{,}05+0{,}95\,\dt\bigr),
\qquad T_{\text{SLS}}^{(1)}=47\,100~\text{lb},
\label{eq:b747-engine-thrust}
\end{equation}
где $\sigma(h)=\rho(h)/\rho_0$~--- относительная плотность атмосферы. Суммарная продольная сила и момент рыскания от асимметричной тяги далее подаются непосредственно в правую часть~\eqref{eq:b747-6dof}:
\begin{equation}
X_{\text{eng}}=\sum_{i=1}^{4}T_i,\qquad
N_{\text{eng}}=-\sum_{i=1}^{4} y_i T_i,\qquad
y=[-71{,}7,\,-35{,}8,\,+35{,}8,\,+71{,}7]~\text{ft}.
\label{eq:b747-engine-moment}
\end{equation}
До отказа распределение тяги симметрично и $N_{\text{eng}}\approx0$. После обнуления $\mu_1$ линия действия оставшейся тяги смещается вправо относительно центра масс: суммарная тяга падает примерно на 25\,\%, а в уравнение рыскания входит ненулевой момент $N_{\text{eng}}$, вызывающий рост $r$, разворот по курсу и связанный крен в сторону отказавшего двигателя. Именно этот скачок оператора ОУ должен быть скомпенсирован каскадом $L_2+L_3$ без передачи контроллеру априорного флага отказа.

\subsection{Параметры контроллера UFTC}

В расчётном эксперименте используется та же реализация \texttt{UFTCController}, но не редуцированная трёхмерная постановка по угловым скоростям, а полный регуляционный вектор:
\begin{equation}
x_{\mathrm{UFTC}} =
\left[
e_V,\ e_h,\ e_{\vartheta},\ e_{\psi},\ e_{\varphi},\ e_{\beta},\ p,\ q,\ r
\right]^\top,
\qquad n_{\text{state}}=9,\quad n_{\text{control}}=4.
\label{eq:uftc-state-b747}
\end{equation}
Здесь $e_V=V-V_0$, $e_h=h-h_0$, $e_{\vartheta}=\vartheta-\vartheta_0$, $e_{\psi}=\psi-\psi_0$, $e_{\varphi}=\varphi-\varphi_0$, $e_{\beta}=\beta-\beta_{\text{ref}}(\hat\mu_1)$; последние три компоненты остаются физическими угловыми скоростями в рад/с. Масштабирование первых шести каналов перед подачей в \UFTC{} задано диагональю
\[
S_x=\mathrm{diag}\left(25~\text{ft/s},\,250~\text{ft},\,3\degr,\,8\degr,\,8\degr,\,2\degr,\,1,\,1,\,1\right).
\]
Модули \IADP{} и \FDD{} работают со всем вектором~\eqref{eq:uftc-state-b747}, а внутренний \AINDI{}-контур берёт из него только индексы $[6,7,8]$, то есть $[p,q,r]$.

Шаг интегрирования $\Delta t=0{,}05$~с, длительность эпизода $T=300$~с, отказ двигателя вводится при $t_f=10$~с. Целевые полосы качества в последнем 20-секундном окне выбраны жёстче текущей настроенной baseline-реализации \ET{}: $\mathrm{RMSE}\,|e_V|\le0{,}050$~ft/s, $\mathrm{RMSE}\,|e_h|\le0{,}01$~ft, $\mathrm{RMSE}\,|e_{\vartheta}|\le0{,}005\degr$, $\mathrm{RMSE}\,|e_{\psi}|\le0{,}010\degr$, $\mathrm{RMSE}\,|e_{\varphi}|\le0{,}010\degr$.

Номинальная балансировка при $h_0=20\,000$~ft и $V_0=674$~ft/s даёт
\[
\alpha_0=\vartheta_0=3{,}603\degr,\qquad
{\dv}_0=-0{,}722\degr,\qquad
{\dt}_0=0{,}5549.
\]
После потери крайнего левого двигателя внешний распределитель равновесного управления использует статическую one-engine-out балансировку:
\[
\beta_{\text{eo}}=-0{,}440\degr,\qquad
{\dv}_{\text{eo}}=-0{,}722\degr,\quad
{\de}_{\text{eo}}=-4{,}403\degr,\quad
{\dn}_{\text{eo}}=-2{,}642\degr,\quad
{\dt}_{\text{eo}}=0{,}7575,
\]
с нормой остатка нелинейных уравнений равновесия $4{,}13\cdot10^{-15}$. В терминах САУ это означает, что канал прямой связи переводит ОУ из симметричного трима в асимметричный трим по измеренной эффективности двигателя, а \UFTC{} работает как онлайн-адаптивная остаточная ОС вокруг этой новой рабочей точки.

Физическая команда формируется суммой трёх слагаемых:
\begin{equation}
u_{\text{phys}} =
u_{\text{trim}}(\hat\mu_1)
+K_x x_{\text{phys}}
+S_u\,u_{\mathrm{UFTC}},
\qquad
S_u=\mathrm{diag}(0,\ 0{,}05\degr,\ 0{,}05\degr,\ 0).
\label{eq:uftc-action-compose}
\end{equation}
Первое слагаемое задаёт равновесную прямую связь, второе является малым полносостояниным законом ОС по $e_V,e_h,e_{\vartheta},e_{\psi},e_{\varphi},e_{\beta},p,q,r$, третье~--- нормированный остаток \UFTC{} в доверительной области. Нулевая остаточная полоса по $\dv$ и $\dt$ выбрана намеренно: продольную энергетику удерживает внешний трим-распределитель с линейной ОС, а онлайн-остаток \UFTC{} демпфирует латерально-направленные и скоростные каналы через элерон и руль направления.

Тёплый старт \IADP{}/\FDD{} задаётся грубой дискретной моделью: $F_{\text{init}}=0{,}99I$ с кинематическими связями $e_{\vartheta}\leftarrow q$, $e_{\psi}\leftarrow r$, $e_{\varphi}\leftarrow p$, а матрица $G_{\text{init}}$ содержит только локальные чувствительности продольной скорости, высоты и угловых скоростей к $[\dv,\de,\dn,\dt]$. Это не идентифицированная модель B-747, а начальное приближение для \RLS{}; в ходе эпизода локальная матрица управления адаптируется онлайн. Параметры \FDD{}: $Q_0=10^{-4}I$, $R_0=10^{-3}I$, \texttt{adapt\_Q=False}, \texttt{adapt\_R=False}, $\mathrm{drift}=6{,}0$, $h_{\text{alarm}}=15{,}0$. Радиусы доверительной области: $\delta_{\text{nom}}=0{,}03$, $\delta_{\text{fault}}=0{,}15$.

\subsection{Сравниваемые подходы}

В качестве референса рассматриваются два контура управления, реализованные в той же программной среде:

\begin{itemize}
\item \emph{Разомкнутый контур.} Рулевые поверхности зафиксированы в исходном балансировочном положении, $\dv = \delta_{\text{в},0}$, $\de = \dn = 0$, $\dt = \delta_{\text{т},0}$. Используется как нижняя граница: показывает, что произошло бы при полном отсутствии регулирования.

\item \textbf{Единый онлайн-актор \ET{}~\cite{mazaev2026etdhp}.} Closed-world архитектура: одна нейронная сеть актор + одна нейронная сеть критик, обученные онлайн методом DHP с событийным триггером. Регуляционное состояние расширено явным признаком отказа $\ell\in\{0,1\}$, что отражает закрытомировое допущение о известности конкретного режима. Модель объекта (одношаговая нейросеть) идентифицируется офлайн на 16\,000 переходов с persistent excitation. Для сопоставления используются текущие baseline-метрики 300-секундного прогона в том же сценарии отказа.
\end{itemize}

\subsection{Результаты UFTC}

После момента отказа $t_f=10$~с обнуление эффективности двигателя~№\,1 создаёт ступенчатый момент рыскания и одновременно уменьшает суммарную продольную тягу. Разомкнутый ОУ быстро уходит из исходного крейсерского режима. В замкнутом контуре сначала измеренная эффективность двигателя переводит $u_{\text{trim}}(\hat\mu_1)$ в асимметричную one-engine-out балансировку; затем линейная ОС подавляет ошибки $e_V,e_h,e_{\vartheta},e_{\psi},e_{\varphi},e_{\beta}$, а остаточный \UFTC{}-контур через \IADP{}+\AINDI{} демпфирует локальные рассогласования модели и угловые скорости. В терминах разработчика САУ алгоритм работает как каскад «перебалансировка рабочей точки + полносостояниная стабилизирующая ОС + адаптивная остаточная компенсация неизвестной динамики».

\begin{table}[H]
\centering
\small
\caption{Пиковые и установившиеся показатели \UFTC{} после отказа двигателя в 300-секундном прогоне.}
\label{tab:uftc-rates}
\begin{tabularx}{\linewidth}{l c c c}
\toprule
Параметр & Пик после отказа & RMSE в последних 20~с & Финальное значение \\
\midrule
$|\Delta V|$, ft/s & $0{,}027$ & $0{,}0010$ & $-0{,}0009$ \\
$|\Delta h|$, ft & $0{,}370$ & $0{,}0012$ & $-0{,}0010$ \\
$|\Delta\vartheta|$, $\degr$ & $0{,}0013$ & $4\cdot10^{-6}$ & $+0{,}000004$ \\
$|\Delta\psi|$, $\degr$ & $0{,}2251$ & $0{,}000620$ & $+0{,}000624$ \\
$|\Delta\varphi|$, $\degr$ & $0{,}8325$ & $0{,}000092$ & $+0{,}000081$ \\
\bottomrule
\end{tabularx}
\end{table}

\paragraph{Анализ диагностического сигнала обнаружителя.} В полном 9-мерном прогоне бинарный сигнал $\text{fault\_present}$ впервые срабатывает в момент $t=10{,}05$~с, то есть через один дискретный шаг после инжекции отказа. Это уже не подпороговый сценарий: ветка «аномалия обнаружена $\to$ расширение доверительной области $\to$ ускорение \RLS{}-переидентификации» реально задействована. Практически важный эффект состоит в том, что \FDD{} не обязан классифицировать тип отказа как «отказ двигателя~№\,1»; ему достаточно зафиксировать статистическую разладку инноваций номинального фильтра, а командная часть САУ использует измеренную эффективность двигателя для перехода на one-engine-out trim.

Качественная картина переходного процесса в замкнутой САУ показана на рис.~\ref{fig:timeline-fault}. До $t=10$~с самолёт находится в симметричном крейсерском триме. В момент отказа исчезает тяга левого внешнего двигателя, скачком меняются продольная сила $X_{\text{eng}}$ и момент рыскания $N_{\text{eng}}$. На следующем шаге \FDD{} выдаёт сигнал \textsc{Тревога}; распределитель переводит прямую связь в асимметричную балансировку, а полносостояниная ОС и \UFTC{}-остаток формируют команды элерона и руля направления, которые гасят $\psi$, $\varphi$, $\beta$ и рейты. В последних 20~с наклоны трендов составляют $+1{,}15\cdot10^{-5}$~ft/s$^2$ по скорости и $+1{,}58\cdot10^{-5}$~ft/s по высоте; дрейф за окно равен $+0{,}00023$~ft/s и $+0{,}00031$~ft соответственно.

\begin{figure}[H]
\centering
\resizebox{\linewidth}{!}{%
\begin{tikzpicture}[>=Stealth,
    lane/.style={rectangle,draw,rounded corners=1pt,minimum height=8mm,align=center,font=\scriptsize},
    sm/.style={font=\scriptsize}]
\def\yplant{0}
\def\yfdd{1.25}
\def\yrls{2.50}
\draw[->,thick] (0,-0.55) -- (12.6,-0.55) node[right,sm] {$t$, с};
\foreach \x/\lab in {0/0,2/10,4.5/11,8/30,12/300} {
  \draw (\x,-0.65) -- (\x,-0.45);
  \node[sm] at (\x,-0.90) {\lab};
}

\node[sm,anchor=east] at (-0.25,\yplant) {ОУ};
\node[sm,anchor=east] at (-0.25,\yfdd) {\FDD{}};
\node[sm,anchor=east,align=right] at (-0.25,\yrls) {RLS};

\draw[fill=green!10,draw=green!50!black] (0,\yplant-0.28) rectangle (2,\yplant+0.28) node[midway,sm] {симм. trim};
\draw[fill=red!18,draw=red!60!black] (2,\yplant-0.28) rectangle (2.25,\yplant+0.28) node[midway,sm] {$t_f$};
\draw[fill=yellow!20,draw=yellow!60!black] (2.25,\yplant-0.28) rectangle (4.5,\yplant+0.28) node[midway,sm,align=center] {engine-out\\trim};
\draw[fill=orange!18,draw=orange!60!black] (4.5,\yplant-0.28) rectangle (8,\yplant+0.28) node[midway,sm,align=center] {затухание\\остатков};
\draw[fill=blue!12,draw=blue!60!black] (8,\yplant-0.28) rectangle (12,\yplant+0.28) node[midway,sm] {удержание режима};

\draw[fill=green!10,draw=green!50!black] (0,\yfdd-0.28) rectangle (2,\yfdd+0.28) node[midway,sm] {норма};
\draw[fill=red!12,draw=red!60!black] (2,\yfdd-0.28) rectangle (2.25,\yfdd+0.28);
\draw[fill=blue!10,draw=blue!60!black] (2.25,\yfdd-0.28) rectangle (12,\yfdd+0.28) node[midway,sm] {severity-модуляция доверительной области};

\draw[fill=gray!10,draw=gray!60] (0,\yrls-0.28) rectangle (2,\yrls+0.28) node[midway,sm,align=center] {ном.\\модель};
\draw[fill=orange!12,draw=orange!60!black] (2,\yrls-0.28) rectangle (12,\yrls+0.28) node[midway,sm] {онлайн-адаптация локальной матрицы управления};

\draw[densely dotted] (2,\yplant+0.45) -- (2,\yrls+0.45);
\draw[densely dotted] (2.25,\yplant+0.45) -- (2.25,\yrls+0.45);
\node[sm,above,align=center] at (3.05,\yrls+0.45) {alarm\\$10{,}05$ с};
\end{tikzpicture}
}
\caption{Временная диаграмма реакции \UFTC{} на отказ двигателя~№\,1 в 300-секундном прогоне. \FDD{} срабатывает в момент $t=10{,}05$~с; после этого прямой канал переводит ОУ на асимметричную балансировку, а \RLS{}+\IADP{}+\AINDI{} подавляют остаточную ошибку регулирования вокруг новой рабочей точки.}
\label{fig:timeline-fault}
\end{figure}

Финальное состояние \UFTC{} в момент $t=300$~с приведено в табл.~\ref{tab:uftc-final}. Разомкнутый контур за тот же горизонт теряет исходный эшелон и уходит в глубокий разворот; \UFTC{} сохраняет исходную скорость и высоту с ошибками порядка $10^{-3}$.

\begin{table}[H]
\centering
\small
\caption{Финальное состояние в момент $t=300$~с по сравнению с разомкнутым контуром.}
\label{tab:uftc-final}
\begin{tabularx}{0.95\linewidth}{l c c}
\toprule
Параметр & Разомкнутый контур & \UFTC{} \\
\midrule
$\Delta V$, ft/s & $+56{,}6$ & $-0{,}0009$ \\
$\Delta h$, ft & $-61\,392$ & $-0{,}0010$ \\
$\Delta\vartheta$, $\degr$ & $-38{,}36$ & $+0{,}000004$ \\
$\Delta\psi$, $\degr$ & $-43{,}04$ & $+0{,}000624$ \\
$\Delta\varphi$, $\degr$ & $-75{,}06$ & $+0{,}000081$ \\
\bottomrule
\end{tabularx}
\end{table}

\begin{figure}[H]
\centering
\includegraphics[width=\linewidth]{figures/uftc_b747_open_vs_closed_vh.png}
\caption{Сравнение разомкнутого контура и \UFTC{} по истинной воздушной скорости и высоте в 300-секундном сценарии отказа двигателя. Разомкнутый ОУ теряет эшелон, тогда как \UFTC{} удерживает $V_0=674$~ft/s и $h_0=20\,000$~ft.}
\label{fig:uftc-open-vs-closed}
\end{figure}

\begin{figure}[H]
\centering
\includegraphics[width=\linewidth]{figures/uftc_b747_diagnostics.png}
\caption{Замкнутый прогон \UFTC{} длительностью 300~с. Сверху вниз: ошибки скорости и высоты; ошибки углов Эйлера; угловые скорости $p,q,r$; команды $\dv,\de,\dn,\dt$; сигналы \FDD{} и \RLS{}.}
\label{fig:uftc-diag}
\end{figure}

\begin{figure}[H]
\centering
\includegraphics[width=0.7\linewidth]{figures/uftc_b747_engine_thrust.png}
\caption{Эффективная тяга двигателей в \UFTC{}-замкнутом контуре. После $t=10$~с тяга двигателя~№\,1 обнуляется; распределитель увеличивает общий уровень $\dt$ и вводит латерально-направленную перебалансировку через $\de,\dn$.}
\label{fig:uftc-thrust}
\end{figure}

\subsection{Сравнение с ET-DHP}

Сравнение \UFTC{} и \ET{}~\cite{mazaev2026etdhp} приводится в табл.~\ref{tab:comparison}. Обе реализации проверяются на одном 300-секундном сценарии отказа внешнего левого двигателя. Отличие постановки сохраняется: \ET{} использует закрытомировой признак отказа $\ell$, а \UFTC{} не расширяет состояние бинарным индикатором режима и обнаруживает статистическую разладку через \FDD{}.

\begin{table}[H]
\centering
\footnotesize
\caption{Сравнительный анализ \UFTC{} и текущей baseline-реализации \ET{} в задаче парирования отказа двигателя B-747.}
\label{tab:comparison}
\begin{tabularx}{\linewidth}{>{\raggedright\arraybackslash}p{0.29\linewidth} >{\raggedright\arraybackslash}X >{\raggedright\arraybackslash}X}
\toprule
Критерий & \UFTC{} & \ET{}~\cite{mazaev2026etdhp} \\
\midrule
Допущение о каталоге отказов & Открытый мир (без каталога) & Закрытый мир (явный признак $\ell$) \\
Расширение состояния признаком отказа & Нет & Да, $\ell\in\{0,1\}$ \\
Офлайн-обучение модели объекта & Не требуется; используется трим-расчёт и онлайн \RLS{} & Требуется (16\,000 PE-переходов, 800 эпох) \\
Онлайн-обучение в эпизоде & Только \RLS{}-идентификация в \IADP{}, без обучения NN & Полный цикл \ET{}-обновления актора и критика \\
Горизонт эпизода & 300~с & 300~с \\
Обнаружитель отказов & Встроен (\FDD{} = Kalman + \CUSUM{}) & Не входит в архитектуру \\
Размерность регуляционного состояния & $9$: ошибки режима + $p,q,r$ & $17$: расширенное состояние \\
Структура управления & Трим-распределитель + полносостояниная ОС + $L_2+L_3+\FDD{}$ & Единый актор + критик \\
Нейросетевые компоненты & Нет & Да \\
\midrule
RMSE $|\Delta V|$ в последних 20~с & $0{,}0010$ ft/s & $0{,}076$ ft/s \\
RMSE $|\Delta h|$ в последних 20~с & $0{,}0012$ ft & $2{,}370$ ft \\
RMSE $\Delta\vartheta,\Delta\psi,\Delta\varphi$ & $4\cdot10^{-6}$; $0{,}000620$; $0{,}000092\degr$ & $0{,}0010$; $0{,}3905$; $0{,}0140\degr$ \\
Итог по заданным полосам качества & \textsc{Pass}; лучше baseline \ET{} & Baseline для сравнения \\
\bottomrule
\end{tabularx}
\end{table}

Полученные числа уточняют интерпретацию первичного прогона. Полносостояниная версия \UFTC{} удерживает исходный крейсерский режим на горизонте 300~с и превосходит текущую baseline-реализацию \ET{} по всем финальным RMSE из табл.~\ref{tab:comparison}. Цена этого результата~--- наличие модельно-зависимого расчёта one-engine-out трима и прямого канала по измеренной эффективности двигателя. Поэтому корректная область превосходства \UFTC{} в данном эксперименте формулируется так: без офлайн-обучения NN и без бинарного признака отказа регулятор достигает более малой установившейся ошибки, если для ОУ доступен статический трим-распределитель отказного режима.

\begin{figure}[H]
\centering
\includegraphics[width=\linewidth]{figures/etdhp_b747_closed_loop.png}
\caption{Замкнутый контур \ET{} длительностью 300~с (источник~--- \cite{mazaev2026etdhp}). Параметры полёта возвращаются к балансировочным значениям с малыми установившимися отклонениями.}
\label{fig:etdhp-closed}
\end{figure}

% ====================================================================
\section{Обсуждение}
\label{sec:disc}
% ====================================================================

\subsection{Применимость UFTC в текущей редакции}

Полносостояниная редакция \UFTC{} применима в задачах, где требуется:
\begin{enumerate}
\item Удержание заданного режима полёта после резкого отказа исполнительного органа при наличии расчётного трим-распределителя.
\item Раннее обнаружение аномалии динамики без поддержания каталога режимов.
\item Перенос между объектами управления без переобучения нейросетей; размерность состояния задаётся разработчиком САУ через выбранные ошибки регулирования.
\item Сохранение работоспособности при отсутствии офлайн-обучения (cold start).
\end{enumerate}

Типичные сценарии~--- внутренний контур стабилизации беспилотных летательных аппаратов в условиях ограниченных вычислительных ресурсов и непредсказуемой эксплуатационной среды; стабилизирующий слой над высокоуровневым автопилотом, реализующим миссию; слой обнаружения, формирующий сигнал тревоги для пилота или системы управления полётом.

\subsection{Ограничения}

Текущая редакция \UFTC{} имеет следующие важные ограничения:

\begin{itemize}
\item \emph{Модельно-зависимый трим-распределитель.} В эксперименте удержание высоты и скорости опирается на расчёт статической one-engine-out балансировки B-747 и на измеренную эффективность двигателя. Без такого канала прямой связи \UFTC{} остаётся только остаточным адаптивным контуром и не гарантирует восстановление продольной энергетики.
\item \emph{Лимиты обнаружения постепенных отказов.} Обнаружитель на основе \CUSUM{} по махаланобисовой дистанции инноваций реагирует на резкие сдвиги среднего и роста дисперсии; постепенный дрейф (например, плавное падение тяги при потере топлива) может не пересекать порог при стандартных настройках. Дальнейшее развитие обнаружителя предусматривает дополнение \CUSUM{} тестом обобщённого отношения правдоподобия (GLR) с фиксированным окном.
\item \emph{Невыполненный формальный анализ устойчивости.} Условия 3-5 раздела~\ref{sec:theory} требуют дальнейшего доказательства; рабочая реализация проверена эмпирически на 56 модульных тестах и интеграционных прогонах, но не имеет полного формального доказательства композитной UUB-границы.
\item \emph{Нет защитного слоя.} HJ-Reachability-фильтр $L_1$ не входит в текущую трёхслойную реализацию, поэтому \UFTC{} в текущей редакции не обеспечивает гарантированного невыхода из envelope-границ при предельных сценариях.
\end{itemize}

\subsection{Программная реализация и воспроизводимость}

\UFTC{} реализован в составе библиотеки \texttt{TensorAeroSpace}~\cite{tensoraerospace}: модуль \texttt{uftc} размещён в пространстве \path{tensoraerospace.agent}. Публичный API: \path{UFTCController(n_state, n_control, nominal_F, nominal_G, config)}. Реализация на pure-NumPy/Gymnasium (без зависимости от GPU и фреймворков глубокого обучения), что обеспечивает применимость на встраиваемых платформах. Покрытие тестами составляет 93\,\% линий кода: 56 модульных тестов и 4 интеграционных теста на нелинейной угловой модели F-16 с шестью пресетами повреждений (отказ руля высоты, потеря руля направления, частичная потеря крыла, отказ двигателя, столкновение с птицей, комбинированный сценарий) и на нелинейной 6-DoF модели Boeing~747 со сценарием отказа крайнего левого двигателя. Реализация поставляется как открытый код под лицензией Apache~2.0; репозиторий~\cite{tensoraerospace}.

% ====================================================================
\section{Заключение}
\label{sec:concl}
% ====================================================================

В работе предложена архитектура унифицированного отказоустойчивого адаптивного контроллера \UFTC{} на основе иерархической композиции инкрементального ADP, активно-адаптивного INDI и параллельного обнаружителя аномалий без априорного каталога режимов отказов. Основные результаты:

\begin{enumerate}
\item Разработана трёхслойная рабочая реализация архитектуры \UFTC{}, в которой средний слой (\IADP{}-обёртка с инициируемым обнаружителем сбросом \RLS{}), внутренний слой (\AINDI{}-обёртка с супер-твист-наблюдателем и непрерывной доверительной областью) и параллельный обнаружитель отказов на основе фильтра Калмана и теста \CUSUM{} образуют связный каскад, управляемый единым скаляром тяжести $s_t\in[0,10]$ и насыщенной величиной $\sigma_t\in[0,1]$.

\item Предложен метод открытомирового обнаружения отказов через тест разладки \CUSUM{} по нормированной махаланобисовой дистанции инноваций единственного номинального фильтра Калмана, не требующий предобучения банка наблюдателей под конкретные режимы. Метод не привязан к фиксированному каталогу отказов и ориентирован на аномалии динамики, лежащие вне заранее заданного набора режимов.

\item Получена корректировка формул адаптации Сейджа-Хузы $Q$ и $R$ с проекцией собственных значений на относительный нижний предел, гарантирующая положительную определённость ковариационных матриц при произвольных начальных условиях. Утверждение~1 даёт формальное доказательство соответствующего свойства.

\item Реализация \UFTC{} в составе библиотеки \texttt{TensorAeroSpace} верифицирована 56 модульными тестами с покрытием 93\,\% и интеграционными прогонами на нелинейных моделях F-16 и Boeing~747. В 300-секундном сценарии отказа крайнего левого двигателя B-747 \FDD{} срабатывает на $t=10{,}05$~с, а \UFTC{} удерживает исходный режим с RMSE в последних 20~с: $0{,}0010$~ft/s по скорости, $0{,}0012$~ft по высоте, $4\cdot10^{-6}\,\degr$, $0{,}000620\degr$, $0{,}000092\degr$ по тангажу, курсу и крену соответственно.
\end{enumerate}

Сравнительный анализ с архитектурой \ET{}~\cite{mazaev2026etdhp} показывает, что в данном сценарии полносостояниная \UFTC{} превосходит текущую baseline-реализацию \ET{} по финальным RMSE без офлайн-обучения нейросетевой модели и без бинарного признака отказа. При этом результат опирается на расчёт one-engine-out трима, поэтому следующая инженерная задача состоит в обобщении трим-распределителя и защитного слоя на более широкий класс отказов.

Перспективы развития работы:
\begin{itemize}
\item Расширение защитного слоя $L_1$ на основе достижимости Гамильтона-Якоби; дополнение обнаружителя \FDD{} тестом обобщённого отношения правдоподобия для постепенных отказов.
\item Реализация внешнего слоя $L_4$ на основе дистрибутивного SAC с фильтром риска CVaR; обобщение удержания продольных параметров на отказы без заранее рассчитанного статического трима.
\item Реализация композитного векторного монитора Ляпунова и формальное доказательство UUB-границы каскада $L_1+L_2+L_3+L_4+\FDD{}+\text{monitor}$.
\item Экспериментальная Monte-Carlo валидация на расширенном наборе пресетов отказов и плантов.
\item Аппаратная валидация (HIL или real flight на UAV X8/quadrotor).
\end{itemize}

% ====================================================================
\section*{Финансирование}
% ====================================================================

Работа выполнена в рамках инициативной научно-исследовательской деятельности автора без целевого финансирования.

\section*{Конфликт интересов}

Автор заявляет об отсутствии конфликта интересов.

% ====================================================================
\section{Vector-comparison UUB-bound каскада UFTC}
\label{sec:uub-lemma}
% ====================================================================

\subsection{Composite Lyapunov-функция}

Введём пять неотрицательных компонент:
\begin{align}
V_{\mathrm{HJ}}(x) &= \max\!\left(0,\, \varepsilon_t - V_\theta(x)\right), \\
V_{\mathrm{INDI}}(\omega,\omega_{\mathrm{ref}}) &= \tfrac12\,\|\omega-\omega_{\mathrm{ref}}\|^2, \\
V_{\mathrm{iADP}}(x_{\mathrm{err}}) &= \tfrac12\,x_{\mathrm{err}}^\top \tilde P\, x_{\mathrm{err}}, \\
V_{\mathrm{DSAC}} &= \max\!\left(0,\, \mathrm{var}(Z) - \overline{\sigma}^2\right), \\
V_{\mathrm{FDD}} &= s_{\mathrm{abrupt}} + s_{\mathrm{gradual}}.
\end{align}
Композитный сигнал $V_{\mathrm{total}}(t) = c^\top v(t)$, где
$v = (V_{\mathrm{HJ}}, V_{\mathrm{INDI}}, V_{\mathrm{iADP}}, V_{\mathrm{DSAC}}, V_{\mathrm{FDD}})^\top$,
$c \in \mathbb{R}_{\ge 0}^5$ — фиксированные веса, $\sum c_i = 1$.

\subsection{Lemma 4.1}

\begin{lemma}[Vector-comparison UUB]
\label{lem:vector-uub}
Пусть для каждой компоненты $V_i$, $i=1,\dots,5$, выполнено
\begin{equation}
\dot V_i(t) \;\le\; -a_i V_i(t) + \sum_{j\ne i} \varepsilon_{ij} V_j(t) + d_i,
\qquad a_i>0,\ \varepsilon_{ij}\ge 0,\ d_i\ge 0,
\end{equation}
почти всюду на траектории системы. Пусть $M = \mathrm{diag}(a) - \varepsilon$ — Hurwitz-Metzler.
Тогда для любого вектора весов $c \in \mathbb{R}^5_{\ge 0}$ с $\|c\|_1=1$
\begin{equation}
V_{\mathrm{total}}(t) \;\le\; \|c\|_1 \cdot \|v(0)\|_\infty \cdot e^{-\lambda_{\min}(M)\, t}
\;+\; \|M^{-1} d\|_c,
\end{equation}
где $\|y\|_c = \sum_i c_i |y_i|$, и в частности
\begin{equation}
\limsup_{t\to\infty} V_{\mathrm{total}}(t) \;\le\; \mu_{\mathrm{uub}} \,\equiv\, \|M^{-1} d\|_c.
\end{equation}
\end{lemma}

\paragraph{Доказательство.}
Применяем принцип векторного сравнения~\cite[\S 9.5]{khalil2002} к $\dot v\le -Mv+d$.
Поскольку $M$ — Hurwitz-Metzler, по результату Перрона-Фробениуса для Metzler-матриц
$M^{-1}\ge 0$ покомпонентно. Получаем покомпонентную оценку
$v(t)\le e^{-Mt}v(0)+M^{-1}d$. Берём $c$-взвешенную сумму. \hfill$\Box$

\subsection{Lemma 4.1' (monitor-augmented).}

При активном Variant-B monitor c macro-actions, эффективные коэффициенты
становятся $a_i' = a_i + \kappa_i$, $\kappa_i\ge 0$, и
$\mu_{\mathrm{uub}}' = \|(M-\mathrm{diag}(\kappa))^{-1} d\|_c \le \mu_{\mathrm{uub}}$.

\subsection{Numerical certificate}

Параметры $(a_i,\varepsilon_{ij},d_i,\kappa_i)$ выводятся явно из
параметров слоёв L1--L4 и FDD каскада UFTC.
Скрипт \texttt{certificate.py} сертифицирует Metzler-структуру,
Hurwitz-спектр и эмпирическую долю траекторий $V_{\mathrm{total}} <
\mu_{\mathrm{uub}}$ на сохранённых rollout'ах по 7 damage-preset'ам;
JSON-отчёт прилагается к артефактам CI.

% ====================================================================
\bibliographystyle{ugost2008}
\bibliography{references}
% ====================================================================

\end{document}
